\subsection{АКТУАЛЬНОСТЬ ТЕМЫ}

В последние годы наблюдается стремительный рост объёмов данных и активное внедрение \gls{rbd} в различные сферы экономики, науки и управления. Однако стандартные средства взаимодействия с \gls{rbd} через язык SQL остаются труднодоступными для пользователей, не обладающих навыками программирования. Решение задачи формирования и выполнения SQL-запросов на основе естественно-языковых запросов пользователя позволит значительно повысить доступность аналитических инструментов и эффективность работы с данными.

Особую актуальность данной работы определяет развитие технологий искусственного интеллекта, появление новых подходов в области обработки естественного языка (\gls{nlp}) и генерации SQL-запросов с использованием больших языковых моделей (\gls{llm}). Интеллектуальные чат-агенты, интегрированные с \gls{rbd}, позволяют не только автоматизировать формирование корректных SQL-запросов, но и обеспечивают наглядную визуализацию результатов анализа, что повышает удобство и скорость принятия решений.

Современные архитектурные решения предусматривают использование таких инструментов, как Python-фреймворки FastAPI и Streamlit, современные СУБД (например, PostgreSQL) и LLM-фреймворки (LangChain, Ollama) для построения полноценных интеллектуальных систем. Разработка чат-агента с поддержкой \gls{rag} технологий предоставляет возможность формирования релевантного контекста для сложных пользовательских запросов, расширяя возможности аналитики и делая её по-настоящему интерактивной.

Таким образом, исследование и реализация интеллектуального чат-агента для работы с реляционными базами данных на естественном языке представляет собой не только научный, но и практический интерес, обеспечивая эффективную обработку и анализ больших объёмов данных в реальном времени.